\documentclass[12pt]{article}
\usepackage{amsmath, amssymb}

\title{CSE400 -- Fundamentals of Probability in Computing\\
Lecture 21--22: Inequalities and Tail Bounds}
\date{}

\begin{document}

\maketitle

\section*{1. Tail and Motivation of Tail Bounds}

\begin{itemize}
    \item Focus on analyzing probabilities of extreme events (tails of distributions).
    \item Objective: Bound probabilities of the form:
    \begin{itemize}
        \item $P(X \geq a)$
        \item $P(|X - \mu| \geq k)$
    \end{itemize}
    \item Motivation:
    \begin{itemize}
        \item Understanding rare events
        \item Application in systems such as predicting worst-case server demand
    \end{itemize}
\end{itemize}

\section*{2. Markov's Inequality}

\subsection*{Definition}
Let $X$ be a non-negative random variable and $a > 0$. Then:
\[
P(X \geq a) \leq \frac{\mathbb{E}[X]}{a}
\]

\subsection*{Assumptions}
\begin{itemize}
    \item $X \geq 0$
    \item $a > 0$
\end{itemize}

\subsection*{Proof}
\[
\mathbb{E}[X] = \int_0^\infty x f(x)\, dx
\]
Split the integral:
\[
\mathbb{E}[X] \geq \int_a^\infty x f(x)\, dx
\]
Since $x \geq a$ for $x \geq a$:
\[
\geq a \int_a^\infty f(x)\, dx
\]
\[
= a \cdot P(X \geq a)
\]
Thus:
\[
P(X \geq a) \leq \frac{\mathbb{E}[X]}{a}
\]

\subsection*{Example}
Example provided in lecture demonstrating application of the bound using expectation.

\section*{3. Chebyshev's Inequality}

\subsection*{Definition}
For a random variable $X$ with mean $\mu$ and variance $\sigma^2$:
\[
P(|X - \mu| \geq k) \leq \frac{\sigma^2}{k^2}
\]

\subsection*{Assumptions}
\begin{itemize}
    \item Mean $\mu$ exists
    \item Variance $\sigma^2$ exists
\end{itemize}

\subsection*{Derivation}
Apply Markov's inequality to:
\[
Y = (X - \mu)^2
\]
Then:
\[
P((X - \mu)^2 \geq k^2) \leq \frac{\mathbb{E}[(X - \mu)^2]}{k^2}
\]
Since:
\[
\mathbb{E}[(X - \mu)^2] = \sigma^2
\]
We obtain:
\[
P(|X - \mu| \geq k) \leq \frac{\sigma^2}{k^2}
\]

\subsection*{Example}
Example in lecture showing deviation bound using variance.

\section*{4. Chernoff Bound and Poisson Tail}

\subsection*{Definition}
Chernoff bound provides an exponential bound on tail probabilities:
\[
P(X \geq a) \leq \inf_{t > 0} \frac{\mathbb{E}[e^{tX}]}{e^{ta}}
\]

\subsection*{Assumptions}
\begin{itemize}
    \item Moment generating function exists
\end{itemize}

\subsection*{Proof (Outline as in lecture)}
\begin{itemize}
    \item Use exponential transformation:
    \[
    P(X \geq a) = P(e^{tX} \geq e^{ta})
    \]
    \item Apply Markov's inequality:
    \[
    \leq \frac{\mathbb{E}[e^{tX}]}{e^{ta}}
    \]
    \item Optimize over $t > 0$
\end{itemize}

\subsection*{Poisson Tail}
\begin{itemize}
    \item Application of Chernoff bound to Poisson random variables.
    \item Produces exponentially decreasing tail probabilities.
\end{itemize}

\subsection*{Example}
Example from lecture applying Chernoff bound to tail probability.

\section*{5. Jensen's Inequality}

\subsection*{Convex Function}
A function $f$ is convex if:
\[
f(\lambda x_1 + (1 - \lambda)x_2) \leq \lambda f(x_1) + (1 - \lambda)f(x_2)
\]

\subsection*{Definition (Jensen's Inequality)}
For a convex function $f$ and random variable $X$:
\[
f(\mathbb{E}[X]) \leq \mathbb{E}[f(X)]
\]

\subsection*{Assumptions}
\begin{itemize}
    \item $f$ is convex
    \item Expectation exists
\end{itemize}

\subsection*{Example}
Example in lecture illustrating convexity and inequality.

\section*{6. Case Study: System-Level Understanding}

\subsection*{Problem}
\begin{itemize}
    \item Predict worst-case demand on a server.
\end{itemize}

\subsection*{Approach}
\begin{itemize}
    \item Model system demand as random variable.
    \item Apply:
    \begin{itemize}
        \item Markov's inequality
        \item Chebyshev's inequality
        \item Chernoff bound
    \end{itemize}
\end{itemize}

\subsection*{Conclusion}
\begin{itemize}
    \item Tail bounds provide guarantees for:
    \begin{itemize}
        \item Worst-case behavior
        \item System reliability
        \item Capacity planning
    \end{itemize}
\end{itemize}

\section*

\end{document}
